\section{Prototype Implementation}\label{sec:impl}

\subsection{Two Tracks}\label{sec:env}
\textbf{Track A -- Solidity on an EVM.} \texttt{IdentityRegistry.sol}, \texttt{ZeroTrustAccess.sol} and \texttt{RevocationGovernance.sol} (Solidity 0.8.24, optimizer 200 runs) implement Section~\ref{sec:arch}: the status machine, Algorithm~1 with anomaly gate and strict mode, \texttt{revokeBatch}/\texttt{suspendBatch}, PEP-authorized root anchoring, and quorum voting. They are compiled with solc-js and executed on the Hardhat in-process EVM. A Mocha/Chai suite of 13 tests passes; besides the registry, replay, revocation and governance cases it covers the industrial additions: the gate denies a high composite score when $A>g$, a tampered anomaly value invalidates the oracle signature, strict resources never step up, batch revocation skips unknown entries and blocks later access, and only authorized PEPs can anchor. Scripts measure gas for every operation. This track gives real gas but, on its own, not multi-node consensus or network latency (Track C adds those), and the device's request signature is verified by the PEP rather than on-chain.

\textbf{Track B -- Python testbed.} A Python implementation of the same state machine uses real Ed25519 (the \texttt{cryptography} library), SHA-256 and a Merkle root. Besides the contract model it contains the \texttt{FastPathPEP}, which verifies the device signature and the oracle attestation (two Ed25519 verifications), applies the per-device counter and tier policy, and logs decisions in Merkle-rooted batches. It is used for scale, statistical, baseline, tier and governance experiments. The host is a 2-vCPU Linux container (Python 3.13, Hardhat 2.22); all seeds derive from 20261004. Track B uses Ed25519 whereas Track A uses secp256k1 ECDSA, so their timings are not comparable.

\textbf{Track C -- Four-validator QBFT network.} The same contracts are deployed on Hyperledger Besu 25.7.0 with four validators (QBFT, 2\,s block period, 4\,s round timeout, chain ID 2026, zero base fee, one JVM of 700\,MB per node, ports 8545--8548) started by \texttt{start\_net.sh} from \texttt{qbftConfigFile.json}. \texttt{besu\_benchmark.js} (ethers v6) measures block interval, sequential call latency, revocation visibility on a second validator, throughput with pre-signed transactions from ten senders, quorum voting and batch recall, and \texttt{besu\_fault.js} stops and restarts validators. Everything, including the load generator, shares the same two-vCPU host, which limits the throughput results (Section~\ref{sec:besu_eval}).

\subsection{Industrial Fleet and Signal Model}\label{sec:devsim}
The synthetic fleet has six classes with proportions sensor 40\%, PLC 15\%, HMI 8\%, robot 7\%, gateway 15\% and legacy 15\%. Native devices draw posture in $[0.75,1]$, anomaly in $[0,0.15]$ and network context in $[0.7,1]$; legacy devices draw posture in $[0.55,0.75]$, anomaly in $[0,0.20]$ and network in $[0.6,0.9]$. To mimic process upsets, a fraction of legitimate requests (5\% by default; 0, 5 and 15\% in the sensitivity study) draws anomaly from $[0.15,0.45]$. Four attack populations are used: \emph{naive compromise} (all signals degraded), \emph{partial mimicry} (posture and network legitimate, $A\in[0.3,0.7]$), \emph{stealthy mimicry} ($A\in[0.15,0.40]$) and \emph{full spoof} ($P=N=1$, $A\in[0.3,0.7]$); each starts from three clean prior requests. The legacy gateway of Section~\ref{sec:legacy} is modelled through the class \texttt{legacy} and its weaker signals; deriving proxy DIDs from the gateway key is a design element we did not implement. In the synthetic experiments the distributions are our assumptions, not field measurements, and the attack and legitimate ranges overlap by construction, so they show how the policy responds to such overlap rather than how well a detector would perform in a plant. The real-data experiment of Section~\ref{sec:real_eval} (\texttt{sim/experiments\_realdata.py}, scikit-learn 1.9) replaces only the behavioural signal.

\subsection{Experiment Scenarios and Reproducibility}\label{sec:scenarios}
\emph{Carried over from the base design:} (S1) scaling over $n\in\{50,\dots,1600\}$; (S2) 500 trials per basic attack; (S3) 30 repeated runs at $n=400$; (S4) baseline revocation exposure for 60 devices at 1 request/s over 600 logical seconds with revocation times from $U(100,500)$\,s; (S5) governance with 200 trials per $N\in\{5,\dots,21\}$ under cooperation, maximal dissent and forced revocation; and a trust-signal ablation. \emph{Industrial additions:} (X1) tiered policy on the mixed fleet, 4{,}000 requests per population; (X2) anomaly-gate sweep; (X3) fast path versus modelled on-chain path, 10 runs of 2{,}000 requests over 300 PLC and robot devices; (X4) recall capacity from measured gas; and (X5) outage scenario; and (R) the TON\_IoT study of Section~\ref{sec:real_eval}, with five seeds for detector quality and one seed for held-out types and the policy runs. The commands are \texttt{python -m sim.benchmark}, \texttt{sim.experiments\_extra}, \texttt{sim.experiments\_governance}, \texttt{sim.experiments\_iiot}, \texttt{sim.experiments\_realdata} (needs the TON\_IoT CSVs in \texttt{data/}; they are not redistributed) and \texttt{hardhat run scripts/measure\_gas.js}, \texttt{measure\_governance\_gas.js}, \texttt{measure\_iiot\_gas.js}; \texttt{paper/gen\_tables.py} regenerates every table from the resulting JSON files.
